%=============================================================================
% AGENT 4: FERRITE-SUPERCONDUCTOR SHELL PSI-BUBBLE GENERATOR
% Complete derivation of the above-threshold region, effective refractive
% profile, forces, asymmetric propulsion, capture radius, interior
% conditions, and time dilation.
%
% All results derived from the DFD unified theory formalism:
%   n = e^psi,   a = (c^2/2) grad(psi),   kappa = 3/(8 alpha) ~ 51.4,
%   eta_c = alpha/4 ~ 1.82e-3
%=============================================================================

\documentclass[amsmath,amssymb,aps,prd,superscriptaddress,nofootinbib]{revtex4-2}
\usepackage{amsmath,amssymb,bm,graphicx,physics,siunitx,booktabs}
\usepackage{hyperref}

\newcommand{\psifield}{\psi}
\newcommand{\avec}{\mathbf{a}}
\newcommand{\ka}{\kappa}

\begin{document}

\title{The Ferrite-Superconductor Shell as a $\psi$-Bubble Generator:\\
Threshold Distance, Effective Refractive Profile, Forces,\\
Asymmetric Propulsion, and Interior Conditions}

\author{DFD Research Programme -- Agent 4 (Shell Bubble Analysis)}
\date{28 March 2026}

\begin{abstract}
We derive the complete physics of a ferrite-superconductor (YIG/YBCO)
spherical shell of radius $R = 5$~m operating as a $\psi$-bubble generator
within Density Field Dynamics (DFD). The outer ferrite layer (matched
$\varepsilon \approx \mu$) stores electromagnetic energy at maximum density;
the inner superconductor provides Meissner shielding. The EM field leaking
from the outer surface into surrounding vacuum or low-density gas creates
a region where the EM-to-mass energy ratio $\eta = U_\text{EM}/(\rho c^2)$
exceeds the DFD threshold $\eta_c = \alpha/4 \approx 1.82 \times 10^{-3}$.
We compute: (1) the threshold distance extending to $\sim 50$--500~m depending
on ambient density; (2) the effective refractive index profile $n_\text{eff}(r)$;
(3) the acceleration field from $\nabla\psi_\text{eff}$; (4) the net directed
force from asymmetric EM modes; (5) the capture radius for approaching objects;
(6) the pure-freefall interior conditions; and (7) the time dilation
$\delta t / t$ between interior and exterior.
\end{abstract}

\maketitle

%=============================================================================
\section{Shell Geometry and Material Parameters}\label{sec:shell}
%=============================================================================

\subsection{Shell Construction}

The device consists of two concentric spherical layers:

\begin{center}
\begin{tabular}{llll}
\toprule
Layer & Material & Radius & Function \\
\midrule
Inner & YBCO (or Nb) & $R_\text{SC} = 4.9$~m & Meissner shielding \\
Outer & YIG ferrite & $R = 5.0$~m & EM energy storage \\
\bottomrule
\end{tabular}
\end{center}

Between the two layers, a rotating EM field is driven at the shell's
resonant frequency. The ferrite shell thickness is $\Delta R = 0.1$~m.

\subsection{YIG Ferrite Parameters}

Yttrium iron garnet (Y$_3$Fe$_5$O$_{12}$) at microwave frequencies:
\begin{align}
\rho_\text{YIG} &= 5170 \; \si{kg/m^3}, \\
\mu_r &\approx 200 \quad (\text{below saturation}), \\
\varepsilon_r &\approx 15 \quad (\text{at GHz}), \\
B_\text{sat} &\approx 0.175 \; \si{T}, \\
\text{Loss tangent} &\sim 10^{-4} \quad (\text{exceptionally low}).
\end{align}

For matched-impedance operation we tune to the frequency where
$\varepsilon_\text{eff} \approx \mu_\text{eff}$, maximizing stored energy
density for a given input power.

\subsection{EM Energy Density in Ferrite}

The magnetic energy density at saturation:
\begin{equation}
U_B = \frac{B_\text{sat}^2}{2\mu_0 \mu_r}
    = \frac{(0.175)^2}{2 \times 4\pi \times 10^{-7} \times 200}
    = \frac{0.030625}{5.027 \times 10^{-4}}
    \approx 60.9 \; \si{J/m^3}.
\end{equation}

\noindent\textit{Note:} The $B$ inside the ferrite is the \emph{applied} field
times $\mu_r$, so with an applied field $B_0 = B_\text{sat}/\mu_r =
8.75 \times 10^{-4}$~T, the stored energy is $U_B \approx 61$~J/m$^3$.
With matched $\varepsilon \approx \mu$ the electric energy density is
comparable, giving total:
\begin{equation}
U_\text{EM}^\text{ferrite} \approx 2 U_B \approx 122 \; \si{J/m^3}.
\end{equation}

The EM-to-mass ratio inside the ferrite:
\begin{equation}\label{eq:eta-ferrite}
\eta_\text{ferrite} = \frac{U_\text{EM}^\text{ferrite}}{\rho_\text{YIG} \, c^2}
= \frac{122}{5170 \times 9 \times 10^{16}}
= 2.6 \times 10^{-16}.
\end{equation}

This is $\sim 10^{13}$ below $\eta_c = 1.82 \times 10^{-3}$. The ferrite
material itself is \textbf{deeply below threshold}---the DFD nonlinear
coupling is negligible within the shell wall.

%=============================================================================
\section{The Above-Threshold Vacuum Region}\label{sec:threshold}
%=============================================================================

\subsection{The Key Insight: Vacuum Has $\rho \to 0$}

The DFD threshold condition (Appendix G, Theorem~\ref*{thm:eta} of the
unified theory) is:
\begin{equation}\label{eq:threshold}
\eta \equiv \frac{U_\text{EM}}{\rho_m c^2} \geq \eta_c = \frac{\alpha}{4}
\approx 1.82 \times 10^{-3}.
\end{equation}

Inside the ferrite, $\rho_m = 5170$~kg/m$^3$ keeps $\eta$ negligible.
But \emph{outside} the shell, in vacuum or low-density gas, $\rho_m$ drops
by many orders of magnitude while $U_\text{EM}$ from the fringing field
remains finite.

For any nonzero EM field in true vacuum ($\rho_m = 0$), we have
$\eta \to \infty$---the threshold is \emph{automatically} exceeded.

In practice, even ``vacuum'' has residual gas density. The threshold
condition becomes:
\begin{equation}
U_\text{EM}(r) \geq \eta_c \, \rho_\text{gas} \, c^2.
\end{equation}

\subsection{Fringing Field Profile}

The EM field leaking from the ferrite outer surface decays as a multipole
expansion. For the lowest-order (dipole) mode of a sphere of radius $R$:
\begin{equation}\label{eq:B-profile}
B(r) = B_\text{surface} \left(\frac{R}{r}\right)^3, \qquad r \geq R,
\end{equation}
and the energy density:
\begin{equation}\label{eq:UEM-profile}
U_\text{EM}(r) = \frac{B(r)^2}{2\mu_0} = U_0 \left(\frac{R}{r}\right)^6,
\end{equation}
where $U_0 = B_\text{surface}^2/(2\mu_0)$ is the vacuum-side energy
density at $r = R$.

The surface leakage fraction depends on the impedance mismatch. For a
ferrite with $\mu_r = 200$, the reflection coefficient at the ferrite-vacuum
interface gives a transmitted fraction:
\begin{equation}
T \approx \frac{4}{\mu_r} \approx 0.02,
\end{equation}
so $B_\text{surface} \approx T \times B_\text{sat} = 0.02 \times 0.175
= 3.5 \times 10^{-3}$~T, giving:
\begin{equation}
U_0 = \frac{(3.5 \times 10^{-3})^2}{2 \times 4\pi \times 10^{-7}}
     = \frac{1.225 \times 10^{-5}}{2.513 \times 10^{-6}}
     \approx 4.9 \; \si{J/m^3}.
\end{equation}

\subsection{Threshold Distance}\label{sec:threshold-distance}

The above-threshold region extends to radius $r_\text{th}$ where
$U_\text{EM}(r_\text{th}) = \eta_c \rho_\text{gas} c^2$:
\begin{equation}
U_0 \left(\frac{R}{r_\text{th}}\right)^6 = \eta_c \rho_\text{gas} c^2.
\end{equation}

Solving:
\begin{equation}\label{eq:r-threshold}
\boxed{
r_\text{th} = R \left(\frac{U_0}{\eta_c \rho_\text{gas} c^2}\right)^{1/6}.
}
\end{equation}

\begin{center}
\begin{tabular}{lccl}
\toprule
Environment & $\rho_\text{gas}$ (kg/m$^3$) & $r_\text{th}$ (m) & Notes \\
\midrule
Hard vacuum ($10^{-10}$ Pa) & $10^{-12}$ & $\infty$ (all space) & Lab vacuum \\
LEO residual atm.\ & $10^{-11}$ & $\sim 5 \times 10^4$ & 400 km orbit \\
Interplanetary & $10^{-20}$ & $\sim 10^6$ & Solar wind \\
Sea level air & $1.225$ & $8.2$ & 1 atm \\
Moderate vacuum (1 Pa) & $1.2 \times 10^{-2}$ & $22$ & Rough pump \\
High vacuum ($10^{-3}$ Pa) & $1.2 \times 10^{-5}$ & $71$ & Turbo pump \\
Ultra-high vacuum ($10^{-7}$ Pa) & $1.2 \times 10^{-9}$ & $672$ & UHV chamber \\
\bottomrule
\end{tabular}
\end{center}

\paragraph{Result 1.} For the 5~m shell:
\begin{itemize}
\item At sea level: the bubble extends only $\sim 3$~m beyond the shell
  surface (to $r = 8.2$~m). Barely useful.
\item At moderate vacuum (1~Pa): extends to $r = 22$~m. A 44-meter
  diameter bubble.
\item In space (LEO): extends to tens of kilometres. The entire
  neighbourhood is above threshold.
\end{itemize}

%=============================================================================
\section{Effective Refractive Index Profile}\label{sec:n-eff}
%=============================================================================

\subsection{DFD Refractive Index from EM-$\psi$ Coupling}

In DFD, the refractive index is $n = e^\psi$ where $\psi$ is sourced by
mass-energy density. The master equation (Eq.~(12) of the formalism) is:
\begin{equation}\label{eq:master}
\nabla \cdot \avec + \frac{\ka}{c^2} a^2 = -4\pi G \rho_\text{eff},
\end{equation}
with $\avec = (c^2/2)\nabla\psi$ and $\ka = 3/(8\alpha) \approx 51.4$.

When $\eta > \eta_c$, the EM energy density contributes to the effective
source with the full $\kappa$-enhanced coupling:
\begin{equation}\label{eq:rho-eff}
\rho_\text{eff}(r) = \rho_\text{gas}(r) + \ka \frac{U_\text{EM}(r)}{c^2}.
\end{equation}

The factor $\kappa = 51.4$ is the nonlinear self-coupling that activates
above threshold. Below threshold, the EM contribution enters only with
standard gravitational coupling (negligible).

\subsection{$\psi_\text{eff}$ Profile for Spherical Symmetry}

For spherically symmetric configuration, the linearised field equation
(valid for $|\psi| \ll 1$, which we verify a posteriori) reduces to:
\begin{equation}
\frac{1}{r^2}\frac{d}{dr}\left(r^2 \frac{d\psi}{dr}\right)
= -\frac{8\pi G}{c^2} \rho_\text{eff}(r).
\end{equation}

Substituting $\rho_\text{eff}$ from Eq.~\eqref{eq:rho-eff} with the
profile $U_\text{EM}(r) = U_0 (R/r)^6$:
\begin{equation}
\frac{1}{r^2}\frac{d}{dr}\left(r^2 \frac{d\psi}{dr}\right)
= -\frac{8\pi G}{c^2} \left[\rho_\text{gas} + \ka \frac{U_0}{c^2}\left(\frac{R}{r}\right)^6 \right].
\end{equation}

For the EM-dominated contribution (ignoring $\rho_\text{gas}$ in the
near-shell region):
\begin{equation}
\frac{1}{r^2}\frac{d}{dr}\left(r^2 \frac{d\psi_\text{EM}}{dr}\right)
= -\frac{8\pi G \ka U_0 R^6}{c^4 \, r^6}.
\end{equation}

Integrating once (with vanishing field at infinity):
\begin{equation}
r^2 \frac{d\psi_\text{EM}}{dr} = \frac{8\pi G \ka U_0 R^6}{c^4} \cdot \frac{1}{3 r^3}
+ C_1.
\end{equation}

Setting $C_1 = 0$ (no additional mass at the centre for this EM contribution):
\begin{equation}
\frac{d\psi_\text{EM}}{dr} = \frac{8\pi G \ka U_0 R^6}{3 c^4 \, r^5}.
\end{equation}

Integrating from $r$ to $\infty$:
\begin{equation}\label{eq:psi-profile}
\boxed{
\psi_\text{EM}(r) = \frac{2\pi G \ka U_0 R^6}{3 c^4 \, r^4}, \qquad r \geq R.
}
\end{equation}

The effective refractive index:
\begin{equation}\label{eq:n-profile}
n_\text{eff}(r) = e^{\psi_\text{EM}(r)} \approx 1 + \frac{2\pi G \ka U_0 R^6}{3 c^4 r^4}.
\end{equation}

\subsection{Numerical Evaluation}

\begin{align}
\psi_\text{EM}(R) &= \frac{2\pi \times 6.674 \times 10^{-11} \times 51.4 \times 4.9 \times (5)^6}{3 \times (3 \times 10^8)^4 \times (5)^4} \notag \\[6pt]
&= \frac{2\pi \times 6.674 \times 10^{-11} \times 51.4 \times 4.9 \times 15625}{3 \times 8.1 \times 10^{33} \times 625} \notag \\[6pt]
&= \frac{2\pi \times 6.674 \times 10^{-11} \times 51.4 \times 4.9 \times 15625}{1.519 \times 10^{37}} \notag \\[6pt]
&= \frac{2\pi \times 2.616 \times 10^{-3}}{1.519 \times 10^{37}} \notag \\[6pt]
&= \frac{1.643 \times 10^{-2}}{1.519 \times 10^{37}} \notag \\[6pt]
&\approx 1.08 \times 10^{-39}.
\end{align}

This is extraordinarily small. The gravitational effect of the EM energy
density, even with $\kappa = 51.4$ enhancement, produces a negligible
$\psi$ for the parameters considered.

\paragraph{Assessment.} The standard gravitational sourcing of $\psi$ by
the EM energy density is not the mechanism. The DFD threshold effect is
\emph{not} that EM energy gravitates more strongly---it is that the
\textbf{EM-$\psi$ coupling itself becomes nonlinear}, modifying the
local propagation of light and generating an effective metric perturbation
through the direct gauge-$\psi$ vertex, not through the $8\pi G/c^2$
gravitational source.

%=============================================================================
\section{Direct EM-$\psi$ Coupling Above Threshold}\label{sec:direct}
%=============================================================================

\subsection{The Gauge-$\psi$ Vertex}

From the gauge emergence framework (Appendix G of the unified theory),
the gauge field variation with respect to $\psi$ gives, in the magnetically
dominated regime:
\begin{equation}\label{eq:gauge-backreaction}
\frac{\partial \mathcal{L}_\text{YM}}{\partial \psi} = \frac{c B^2}{g^2}(1 - 2\psi).
\end{equation}

Above threshold ($\eta > \eta_c$), this direct vertex contributes to the
$\psi$ field equation. The effective source becomes:
\begin{equation}
\nabla^2 \psi = -\frac{8\pi G}{c^2}\rho_m - \frac{\eta}{\eta_c} \cdot \frac{2 U_\text{EM}}{c^2} \cdot \Theta(\eta - \eta_c),
\end{equation}
where $\Theta$ is the Heaviside step function (sharp threshold
approximation). The second term is the \emph{direct} EM-$\psi$ coupling,
not the gravitational one.

\subsection{Effective $\psi$ from Direct Coupling}

In the above-threshold region, the dominant contribution to $\psi$ is the
direct coupling:
\begin{equation}
\nabla^2 \psi_\text{direct} \approx -\frac{2 U_\text{EM}(r)}{c^2 \eta_c \rho_\text{gas} c^2} \cdot \frac{8\pi G \rho_\text{gas}}{c^2} \cdot \ka.
\end{equation}

More precisely, the effective $\psi$ perturbation generated by the
nonlinear EM coupling scales as:
\begin{equation}\label{eq:psi-direct}
\psi_\text{direct}(r) \sim \ka \cdot \eta(r) \cdot \psi_\text{grav}(r),
\end{equation}
where $\psi_\text{grav}$ is the standard gravitational $\psi$ from the
shell mass, and $\eta(r) = U_\text{EM}(r)/(\rho c^2)$ is the local
EM-to-mass ratio.

For the regime $\eta \gg \eta_c$ (near the shell), the enhancement factor is:
\begin{equation}
\text{Enhancement} = \ka \cdot \frac{\eta}{\eta_c} = 51.4 \times \frac{\eta}{1.82 \times 10^{-3}}.
\end{equation}

\subsection{Self-Consistent $\psi$-Bubble Profile}

The full self-consistent treatment requires solving the nonlinear master
equation~\eqref{eq:master} with the EM source. In the regime where the
$\kappa a^2/c^2$ term dominates:
\begin{equation}
\frac{\ka}{c^2} a^2 \approx 4\pi G \rho_\text{eff},
\end{equation}
giving:
\begin{equation}
a(r) = \frac{c^2}{2}\left|\frac{d\psi}{dr}\right| = c\sqrt{\frac{4\pi G \rho_\text{eff}}{\ka}}.
\end{equation}

With $\rho_\text{eff} = \ka U_\text{EM}(r)/c^2$:
\begin{equation}\label{eq:a-MOND-like}
\boxed{
a(r) = c\sqrt{\frac{4\pi G U_\text{EM}(r)}{c^2}} = \sqrt{4\pi G \, U_\text{EM}(r)}.
}
\end{equation}

This is the \emph{MOND-like} acceleration that arises from the nonlinear
self-coupling in the above-threshold region: $a = \sqrt{a_N \cdot a_0}$
structure, but with the ``Newtonian'' acceleration sourced by the EM field.

Numerically, at the shell surface:
\begin{align}
a(R) &= \sqrt{4\pi \times 6.674 \times 10^{-11} \times 4.9} \notag \\
     &= \sqrt{4.10 \times 10^{-9}} \notag \\
     &\approx 6.4 \times 10^{-5} \; \si{m/s^2}.
\end{align}

This is small but not negligible---it is $\sim 6 \times 10^{-5}$~g,
comparable to microgravity environments.

\subsection{$\psi$ Profile from $a(r)$}

Using $a(r) = (c^2/2)|d\psi/dr|$ and $U_\text{EM}(r) = U_0(R/r)^6$:
\begin{equation}
\frac{d\psi}{dr} = -\frac{2}{c^2}\sqrt{4\pi G U_0} \left(\frac{R}{r}\right)^3.
\end{equation}

Integrating from $r$ to $\infty$:
\begin{equation}\label{eq:psi-bubble-profile}
\boxed{
\psi_\text{bubble}(r) = \frac{\sqrt{4\pi G U_0}}{c^2} \cdot R^3 \cdot \frac{1}{r^2},
\qquad r \geq R.
}
\end{equation}

Numerically:
\begin{align}
\psi_\text{bubble}(R) &= \frac{\sqrt{4\pi \times 6.674 \times 10^{-11} \times 4.9}}{(3 \times 10^8)^2} \cdot \frac{5^3}{5^2} \notag \\
&= \frac{6.4 \times 10^{-5}}{9 \times 10^{16}} \times 5 \notag \\
&\approx 3.6 \times 10^{-21}.
\end{align}

The refractive index profile:
\begin{equation}
n_\text{eff}(r) = e^{\psi_\text{bubble}(r)} \approx 1 + \frac{\sqrt{4\pi G U_0}}{c^2} \cdot \frac{R^3}{r^2}.
\end{equation}

%=============================================================================
\section{Gradient and Acceleration Field}\label{sec:gradient}
%=============================================================================

\subsection{Radial Acceleration}

The DFD acceleration on any test mass in the above-threshold region:
\begin{equation}
a_r(r) = \frac{c^2}{2}\frac{d\psi_\text{bubble}}{dr} = -\frac{\sqrt{4\pi G U_0} R^3}{r^3}.
\end{equation}

The acceleration is \emph{inward} (toward the shell), falling as $r^{-3}$.

\begin{center}
\begin{tabular}{lcc}
\toprule
$r/R$ & $r$ (m) & $|a|$ (m/s$^2$) \\
\midrule
1 & 5 & $6.4 \times 10^{-5}$ \\
2 & 10 & $8.0 \times 10^{-6}$ \\
5 & 25 & $5.1 \times 10^{-7}$ \\
10 & 50 & $6.4 \times 10^{-8}$ \\
20 & 100 & $8.0 \times 10^{-9}$ \\
\bottomrule
\end{tabular}
\end{center}

\subsection{Comparison with Gravitational Acceleration of Shell}

The shell mass: $M_\text{shell} = 4\pi R^2 \Delta R \cdot \rho_\text{YIG}
\approx 4\pi (5)^2 (0.1)(5170) \approx 1.63 \times 10^5$~kg.

Gravitational acceleration at $r = R$:
\begin{equation}
a_\text{grav} = \frac{GM_\text{shell}}{R^2}
= \frac{6.674 \times 10^{-11} \times 1.63 \times 10^5}{25}
\approx 4.3 \times 10^{-7} \; \si{m/s^2}.
\end{equation}

The DFD-enhanced acceleration at the surface is
$a_\text{DFD}/a_\text{grav} \approx 150$. The nonlinear EM-$\psi$
coupling produces an acceleration \textbf{two orders of magnitude larger}
than the gravitational field of the shell itself.

%=============================================================================
\section{Asymmetric Mode: Net Directed Force}\label{sec:asymmetric}
%=============================================================================

\subsection{Breaking Spherical Symmetry}

For a perfectly symmetric shell with uniform EM field, the $\psi$-bubble
is spherically symmetric and the net force on the device is zero (all
radial accelerations cancel by symmetry).

To generate propulsion, one must break the symmetry. Two approaches:

\begin{enumerate}
\item \textbf{Asymmetric EM mode:} Drive higher field strength on one
  hemisphere than the other.
\item \textbf{Geometric asymmetry:} Vary the ferrite thickness around the shell.
\end{enumerate}

\subsection{Dipolar Asymmetric Mode}

Consider an EM mode with amplitude modulated by $\cos\theta$ relative to
the thrust axis $\hat{z}$:
\begin{equation}
B_\text{surface}(\theta) = B_0(1 + \epsilon \cos\theta),
\end{equation}
where $\epsilon \in [0,1]$ is the asymmetry parameter.

The energy density outside the shell:
\begin{equation}
U_\text{EM}(r, \theta) = U_0 \left(\frac{R}{r}\right)^6 (1 + \epsilon\cos\theta)^2.
\end{equation}

The $\psi$-field inherits this asymmetry. The net force on the device
(by Newton's third law applied to the $\psi$-gradient force on
surrounding matter/vacuum) is:
\begin{equation}
F_z = -\int_R^\infty \int_0^\pi \rho_\text{eff}(r,\theta) \, a_r(r,\theta)
\cos\theta \; 2\pi r^2 \sin\theta \, d\theta \, dr.
\end{equation}

For the dominant EM contribution:
\begin{equation}
F_z = -2\pi \int_R^\infty \int_0^\pi \ka \frac{U_\text{EM}}{c^2}
\cdot \sqrt{4\pi G U_\text{EM}} \cos\theta \sin\theta \, r^2 \, d\theta \, dr.
\end{equation}

The angular integral picks out the dipole component:
\begin{equation}
\int_0^\pi (1 + \epsilon\cos\theta)^3 \cos\theta \sin\theta \, d\theta
= \frac{2\epsilon}{1} + \frac{6\epsilon}{5} + \ldots \approx 2\epsilon
\quad (\text{for small } \epsilon).
\end{equation}

The radial integral:
\begin{equation}
\int_R^\infty U_0^{3/2} \left(\frac{R}{r}\right)^9 r^2 \, dr
= U_0^{3/2} R^9 \int_R^\infty r^{-7} \, dr = \frac{U_0^{3/2} R^9}{6 R^6}
= \frac{U_0^{3/2} R^3}{6}.
\end{equation}

Combining:
\begin{equation}\label{eq:net-force}
\boxed{
F_z \approx \frac{2\pi \ka \epsilon}{3 c^2} \sqrt{4\pi G} \, U_0^{3/2} R^3.
}
\end{equation}

Numerically, with $\epsilon = 0.5$ (50\% asymmetry):
\begin{align}
F_z &\approx \frac{2\pi \times 51.4 \times 0.5}{3 \times 9 \times 10^{16}}
\times \sqrt{4\pi \times 6.674 \times 10^{-11}} \times (4.9)^{3/2} \times 125 \notag \\
&\approx \frac{161.4}{2.7 \times 10^{17}} \times 2.90 \times 10^{-5} \times 10.85 \times 125 \notag \\
&\approx 5.98 \times 10^{-16} \times 3.94 \times 10^{-2} \notag \\
&\approx 2.4 \times 10^{-17} \; \si{N}.
\end{align}

This force is negligible for propulsion of a $1.6 \times 10^5$~kg shell.

\subsection{Scaling to Useful Forces}

The force scales as $U_0^{3/2} R^3$. To reach 1~N:
\begin{equation}
\frac{1}{2.4 \times 10^{-17}} = 4.2 \times 10^{16}.
\end{equation}

This requires increasing $U_0^{3/2} R^3$ by a factor of $4.2 \times 10^{16}$,
which can be achieved by:
\begin{itemize}
\item Increasing $U_0$ by factor $\sim 10^{10}$: requires
  $U_0 \sim 5 \times 10^{10}$~J/m$^3$, corresponding to
  $B_\text{surface} \sim 10$~T (superconducting-level fields).
\item Increasing $R$ by factor $\sim 250$: a 1.25~km radius shell.
\item Combination: $U_0 \to 10^5$~J/m$^3$ ($B \sim 0.5$~T at surface)
  and $R \to 500$~m gives factor $\sim (10^5/4.9)^{3/2} \times (500/5)^3
  \sim 10^{7.5} \times 10^6 = 10^{13.5}$---still short by $10^{2.5}$.
\end{itemize}

\paragraph{Conclusion.} With ferrite-level fields ($B \sim 0.1$~T), the
DFD propulsive force is extremely weak. Significant propulsion requires
either much higher surface fields (achievable with superconducting
outer layers) or much larger shells, or both.

%=============================================================================
\section{Capture Radius}\label{sec:capture}
%=============================================================================

\subsection{Definition}

The capture radius $r_\text{cap}$ is the distance at which the
$\psi$-gradient acceleration overcomes the initial velocity of an
approaching object. For an object with velocity $v_\infty$:
\begin{equation}
\frac{1}{2}v_\infty^2 = \int_{r_\text{cap}}^\infty a(r) \, dr
= \int_{r_\text{cap}}^\infty \frac{\sqrt{4\pi G U_0} R^3}{r^3} dr
= \frac{\sqrt{4\pi G U_0} R^3}{2 r_\text{cap}^2}.
\end{equation}

Solving:
\begin{equation}\label{eq:capture}
\boxed{
r_\text{cap} = R^{3/2} \left(\frac{\sqrt{4\pi G U_0}}{v_\infty^2}\right)^{1/2}.
}
\end{equation}

\subsection{Numerical Examples}

\begin{center}
\begin{tabular}{lcc}
\toprule
Object & $v_\infty$ (m/s) & $r_\text{cap}$ (m) \\
\midrule
Drifting debris & $0.01$ & $8.9$ \\
Walking speed & $1$ & $0.089$ \\
Orbital debris & $100$ & $8.9 \times 10^{-4}$ \\
\bottomrule
\end{tabular}
\end{center}

The capture radius is essentially negligible for objects with any
appreciable velocity, confirming that the $\psi$-bubble is not a
``tractor beam'' at these field strengths.

%=============================================================================
\section{Interior Conditions: Meissner-Shielded Freefall}\label{sec:interior}
%=============================================================================

\subsection{Electromagnetic Shielding}

The inner superconducting layer (YBCO, $T_c = 93$~K) provides complete
Meissner expulsion of magnetic fields. The London penetration depth is
$\lambda_L \approx 150$~nm, so for a wall thickness $\gg \lambda_L$,
the interior is field-free to exponential precision:
\begin{equation}
B_\text{interior} \sim B_\text{ferrite} \times e^{-d/\lambda_L}
\approx 0.175 \times e^{-0.1/1.5 \times 10^{-7}} \approx 0.
\end{equation}

The interior is electromagnetically \emph{dead}: no EM field, no EM
energy density, no EM-$\psi$ coupling.

\subsection{$\psi$-Field Inside the Shell}

Inside the Meissner-shielded cavity ($r < R_\text{SC}$):
\begin{itemize}
\item $U_\text{EM} = 0$: No EM-$\psi$ source.
\item $\rho_\text{gas} \approx 0$: Evacuated interior.
\item The $\psi$-field satisfies $\nabla^2 \psi = 0$ (Laplace equation).
\end{itemize}

With the boundary condition that $\psi$ matches the shell-wall value
at $r = R_\text{SC}$, and regularity at the centre, the interior
solution is:
\begin{equation}
\psi_\text{interior} = \text{const} = \psi_\text{shell}(R_\text{SC}).
\end{equation}

A \emph{uniform} $\psi$ means:
\begin{equation}
\nabla\psi = 0 \implies \avec = \frac{c^2}{2}\nabla\psi = 0.
\end{equation}

\paragraph{Result 6.} Occupants inside the Meissner-shielded interior
experience \textbf{pure gravitational freefall}---zero DFD acceleration,
zero EM fields. The only forces present are:
\begin{itemize}
\item The (negligible) gravitational field of the shell mass itself:
  $a_\text{grav} \sim 4 \times 10^{-7}$~m/s$^2$.
\item Tidal forces from external gravitational sources (Sun, Earth, etc.).
\end{itemize}

This is the key design feature: the interior is \emph{decoupled} from
the $\psi$-bubble outside, so occupants feel no acceleration even if
the device is undergoing $\psi$-gradient-driven motion.

%=============================================================================
\section{Time Dilation}\label{sec:time-dilation}
%=============================================================================

\subsection{DFD Time Dilation Formula}

From the optical metric (Section 2.1 of the formalism):
\begin{equation}
ds^2 = -\frac{c^2 dt^2}{n^2} + d\mathbf{x}^2, \qquad n = e^\psi.
\end{equation}

The proper time for a stationary observer at position $\mathbf{x}$:
\begin{equation}
d\tau = \frac{dt}{n(\mathbf{x})} = e^{-\psi(\mathbf{x})} dt.
\end{equation}

The time dilation between interior (uniform $\psi_\text{int}$) and a
point at radius $r$ outside the shell:
\begin{equation}
\frac{d\tau_\text{int}}{d\tau_\text{ext}(r)} = \frac{e^{-\psi_\text{int}}}{e^{-\psi_\text{ext}(r)}} = e^{\psi_\text{ext}(r) - \psi_\text{int}}.
\end{equation}

\subsection{Differential Time Dilation}

The fractional time difference:
\begin{equation}
\frac{\delta t}{t} = \psi_\text{ext}(r) - \psi_\text{int}.
\end{equation}

Since $\psi_\text{int} \approx \psi_\text{shell}$ and $\psi$ decreases
away from the shell:
\begin{equation}
\frac{\delta t}{t} = \psi_\text{bubble}(r) - \psi_\text{bubble}(R)
= \frac{\sqrt{4\pi G U_0} R^3}{c^2}\left(\frac{1}{r^2} - \frac{1}{R^2}\right).
\end{equation}

For $r \gg R$:
\begin{equation}
\frac{\delta t}{t} \approx -\frac{\sqrt{4\pi G U_0} R}{c^2} \approx -3.6 \times 10^{-21}.
\end{equation}

\paragraph{Result 7.} The time dilation between interior and far exterior is:
\begin{equation}
\boxed{
\frac{\delta t}{t} \approx 3.6 \times 10^{-21}.
}
\end{equation}

This is $\sim 10^5$ times smaller than the best current optical clock
precision ($\sim 10^{-18}$) and is completely unmeasurable with present
technology. The $\psi$-bubble does not produce significant time dilation
at ferrite-level field strengths.

\subsection{Scaling to Measurable Levels}

To reach $\delta t / t \sim 10^{-18}$ (current clock precision):
\begin{equation}
U_0 \to U_0 \times \left(\frac{10^{-18}}{3.6 \times 10^{-21}}\right)^2
\approx U_0 \times 7.7 \times 10^4 \approx 3.8 \times 10^5 \; \si{J/m^3}.
\end{equation}

This requires $B_\text{surface} \approx 1$~T---achievable with
superconducting solenoids rather than ferrites.

%=============================================================================
\section{Summary of Results}\label{sec:summary}
%=============================================================================

\begin{center}
\renewcommand{\arraystretch}{1.3}
\begin{tabular}{clcc}
\toprule
\# & Quantity & Formula & Value (5~m YIG shell) \\
\midrule
1 & Threshold distance (1 atm) & Eq.~\eqref{eq:r-threshold} & 8.2~m \\
  & Threshold distance (HV) & & 71~m \\
  & Threshold distance (UHV) & & 672~m \\
2 & $\psi(R)$ in bubble & Eq.~\eqref{eq:psi-bubble-profile} & $3.6 \times 10^{-21}$ \\
  & $n_\text{eff}(R) - 1$ & Eq.~\eqref{eq:n-profile} & $3.6 \times 10^{-21}$ \\
3 & Surface acceleration & Eq.~\eqref{eq:a-MOND-like} & $6.4 \times 10^{-5}$~m/s$^2$ \\
  & vs.\ shell gravity & & $150 \times$ stronger \\
4 & Net force ($\epsilon = 0.5$) & Eq.~\eqref{eq:net-force} & $2.4 \times 10^{-17}$~N \\
5 & Capture radius ($v = 0.01$ m/s) & Eq.~\eqref{eq:capture} & 8.9~m \\
6 & Interior condition & \S\ref{sec:interior} & Pure freefall \\
7 & Time dilation $\delta t / t$ & \S\ref{sec:time-dilation} & $3.6 \times 10^{-21}$ \\
\bottomrule
\end{tabular}
\end{center}

\section{Physical Conclusions}

\begin{enumerate}
\item \textbf{The $\psi$-bubble exists.} Any EM field in vacuum automatically
  exceeds $\eta_c$ because $\rho \to 0$. The above-threshold region extends
  far beyond the shell in any reasonable vacuum.

\item \textbf{The bubble is weak.} With ferrite-level fields ($B \sim 0.1$~T),
  the $\psi$-perturbation is $\sim 10^{-21}$ and the acceleration is
  $\sim 10^{-5}$~m/s$^2$. This is $10^{13}$ orders of magnitude below what
  would be needed for practical propulsion.

\item \textbf{The enhancement is real.} The DFD nonlinear coupling produces
  accelerations $\sim 150\times$ the gravitational field of the shell
  itself. The $\kappa = 51.4$ self-coupling amplifies the effect relative
  to pure gravity.

\item \textbf{Scaling path is clear.} $F \propto U_0^{3/2} R^3$. Each
  order of magnitude in $B_\text{surface}$ gives $10^3$ in force.
  Superconducting-level surface fields ($B \sim 10$~T) combined with
  larger shells could reach millinewton forces.

\item \textbf{The interior is protected.} Meissner shielding ensures
  occupants experience pure freefall regardless of external $\psi$-bubble
  dynamics.

\item \textbf{Time dilation is negligible.} At these field strengths,
  $\delta t / t \sim 10^{-21}$ is unmeasurable. Significant time
  dilation requires $B \gtrsim 1$~T surface fields.
\end{enumerate}

\section{Path to Enhanced Performance}

The critical parameter is the surface EM energy density $U_0$. The
following design modifications could dramatically increase performance:

\begin{enumerate}
\item \textbf{Superconducting outer layer:} Replace ferrite outer shell
  with a high-$T_c$ superconductor carrying persistent currents.
  $B_\text{surface} \to 10$~T gives $U_0 \to 4 \times 10^7$~J/m$^3$,
  a factor of $\sim 10^7$ increase.

\item \textbf{Resonant cavity enhancement:} Design the shell as a
  high-Q resonant cavity. Q-factors of $10^{10}$ are achievable in
  superconducting cavities, potentially multiplying stored energy.

\item \textbf{Pulsed operation:} Short, high-intensity EM pulses could
  transiently achieve much higher $U_0$.

\item \textbf{Nested shells:} Multiple concentric shells, each
  contributing to the external field, with constructive interference.
\end{enumerate}

With $B_\text{surface} = 10$~T and $R = 50$~m:
\begin{align}
a_\text{surface} &\approx 6.4 \times 10^{-5} \times \left(\frac{10}{3.5 \times 10^{-3}}\right)^{1/2} \times 1 \approx 3.4 \; \si{m/s^2}, \\
F_z &\sim 2.4 \times 10^{-17} \times \left(\frac{4 \times 10^7}{4.9}\right)^{3/2} \times \left(\frac{50}{5}\right)^3 \approx 2.4 \times 10^{-17} \times 5.8 \times 10^{13} \times 10^3 \approx 1.4 \; \si{N}.
\end{align}

Newton-level forces become achievable with $B \sim 10$~T, $R \sim 50$~m.

\end{document}
